% 附录 B 分位数表与内容对照
\chapter{分位数表与内容对照}
\label{app:tables}

表中 $\nu=\infty$ 一行即标准正态分布的分位数 $z_q$。常用：双侧 90\%、95\%、99\% 区间分别用 $1.645$、$1.960$、$2.576$；单侧 5\% 检验用 $1.645$。

\noindent
\begin{minipage}[t]{0.56\linewidth}
\subsection*{B.1\quad \texorpdfstring{$t$}{t} 分布的分位数 \texorpdfstring{$t_{\nu,q}$}{t(ν,q)}}
\footnotesize\linespread{1.0}\selectfont\renewcommand{\arraystretch}{1.08}
\begin{tabular}{c*{5}{c}}
\toprule
$\nu$ & $q=0.90$ & $0.95$ & $0.975$ & $0.99$ & $0.995$\\
\midrule
1 & 3.078 & 6.314 & 12.706 & 31.821 & 63.657 \\
2 & 1.886 & 2.920 & 4.303 & 6.965 & 9.925 \\
3 & 1.638 & 2.353 & 3.182 & 4.541 & 5.841 \\
4 & 1.533 & 2.132 & 2.776 & 3.747 & 4.604 \\
5 & 1.476 & 2.015 & 2.571 & 3.365 & 4.032 \\
6 & 1.440 & 1.943 & 2.447 & 3.143 & 3.707 \\
7 & 1.415 & 1.895 & 2.365 & 2.998 & 3.499 \\
8 & 1.397 & 1.860 & 2.306 & 2.896 & 3.355 \\
9 & 1.383 & 1.833 & 2.262 & 2.821 & 3.250 \\
10 & 1.372 & 1.812 & 2.228 & 2.764 & 3.169 \\
12 & 1.356 & 1.782 & 2.179 & 2.681 & 3.055 \\
15 & 1.341 & 1.753 & 2.131 & 2.602 & 2.947 \\
20 & 1.325 & 1.725 & 2.086 & 2.528 & 2.845 \\
25 & 1.316 & 1.708 & 2.060 & 2.485 & 2.787 \\
30 & 1.310 & 1.697 & 2.042 & 2.457 & 2.750 \\
40 & 1.303 & 1.684 & 2.021 & 2.423 & 2.704 \\
60 & 1.296 & 1.671 & 2.000 & 2.390 & 2.660 \\
120 & 1.289 & 1.658 & 1.980 & 2.358 & 2.617 \\
$\infty$ & 1.282 & 1.645 & 1.960 & 2.326 & 2.576 \\
\bottomrule
\end{tabular}
\end{minipage}\hfill
\begin{minipage}[t]{0.4\linewidth}
\subsection*{B.2\quad 卡方分布的分位数 \texorpdfstring{$\chi^2_{\nu,q}$}{chi2(ν,q)}}
\footnotesize\linespread{1.0}\selectfont\renewcommand{\arraystretch}{1.08}
\begin{tabular}{c*{4}{c}}
\toprule
$\nu$ & $q=0.025$ & $0.05$ & $0.95$ & $0.975$\\
\midrule
1 & 0.001 & 0.004 & 3.841 & 5.024 \\
2 & 0.051 & 0.103 & 5.991 & 7.378 \\
3 & 0.216 & 0.352 & 7.815 & 9.348 \\
4 & 0.484 & 0.711 & 9.488 & 11.143 \\
5 & 0.831 & 1.145 & 11.070 & 12.833 \\
6 & 1.237 & 1.635 & 12.592 & 14.449 \\
7 & 1.690 & 2.167 & 14.067 & 16.013 \\
8 & 2.180 & 2.733 & 15.507 & 17.535 \\
9 & 2.700 & 3.325 & 16.919 & 19.023 \\
10 & 3.247 & 3.940 & 18.307 & 20.483 \\
15 & 6.262 & 7.261 & 24.996 & 27.488 \\
20 & 9.591 & 10.851 & 31.410 & 34.170 \\
30 & 16.791 & 18.493 & 43.773 & 46.979 \\
\bottomrule
\end{tabular}
\end{minipage}

\clearpage
\section*{B.3\quad 原讲义内容的去向}

\begin{center}\small
\begin{tabularx}{\linewidth}{>{\raggedright\arraybackslash}p{4.6cm} >{\raggedright\arraybackslash}X}
\toprule
原讲义 & 本讲义 \\
\midrule
1.1–1.3 样本、统计量、抽样分布、样本方差 & 第 1 章 1.1–1.4，补充睡眠数据与图 \\
1.4 简单随机抽样与分层抽样 & 1.5–1.6（协方差已在第 0 章讲过），新增 Neyman 分配；例 1.1、1.2 移入 \starmark{} 深入讨论 \\
1.5 只保存部分数据 & 1.7 \starmark{} 深入讨论（精简） \\
2.1–2.3 矩估计、可识别性、似然 & 第 3 章 3.1–3.4 \\
2.4、2.6 偏差、MSE、相合性 & 第 2 章，新增均匀分布估计量的完整比较与图 \\
2.5 不完整记录与随机停止 & 3.6 \starmark{} 截尾数据、记录方式与 Fisher 信息 \\
2.6 渐近正态、Delta 方法 & 第 0 章 0.6–0.7 先讲工具，第 3 章新增 Fisher 信息与 Cramér–Rao 下界 \\
2.7、第 3 章 无偏构造与估计量比较 & 第 2 章 2.4–2.5，2.7 \starmark{}（对称化、Rao–Blackwell 思想） \\
第 4 章 检验的概率逻辑 & 第 4 章，新增置信区间的覆盖、$t$ 分布推导、Wilson 区间、功效与样本量、p 值误读；随机化检验移入 \starmark{} \\
第 5 章 合法对称性与精确检验 & 第 5 章 5.6 置换检验与符号翻转（以睡眠数据为例），其余（游程等）从略 \\
第 6 章 单样本与双样本 & 第 4 章 4.2、第 5 章；新增 \textbf{Welch 检验}、模拟比较、比例比较案例、QQ 图与敏感性分析 \\
第 7 章 一元线性回归 & 第 6 章；新增几何图示、置信带与预测带、残差图、杠杆点图、Forbes 完整案例、辛普森悖论 \\
53 道章末习题与各节例题 & 选取 13 道作为 C 级挑战题，另新编 A、B 级习题，全部附解答 \\
\bottomrule
\end{tabularx}
\end{center}
